% Prodotto vettoriale c = a x b: perpendicolare al piano di a e b, verso dato
% dalla regola della mano destra; il modulo è l'area del parallelogramma.
\begin{tikzpicture}[>=Stealth, font=\small]
  % piano di a e b, in prospettiva
  \fill[formula!8] (-0.6,-0.5) -- (4.6,-0.5) -- (5.8,1.4) -- (0.6,1.4) -- cycle;
  \node[formula!80, anchor=north east] at (4.6,-0.55) {piano di $\vec a$ e $\vec b$};
  % parallelogramma
  \fill[esempio!25] (0,0) -- (3.2,0) -- (4.6,1.0) -- (1.4,1.0) -- cycle;
  \draw[dashed, gray] (3.2,0) -- (4.6,1.0) -- (1.4,1.0);
  \node[esempio!80!black] at (2.35,0.45) {area $= ab\sin\theta$};
  \draw[->, very thick, formula] (0,0) -- (3.2,0) node[below] {$\vec a$};
  \draw[->, very thick, medica] (0,0) -- (1.4,1.0) node[above left] {$\vec b$};
  \draw[esempio, thick] (0.6,0) arc[start angle=0, end angle=35.5, radius=0.6];
  \node[esempio] at (17:0.85) {$\theta$};
  % c perpendicolare
  \draw[->, very thick, attenzione] (0,0) -- (0,2.2) node[above] {$\vec c = \vec a \times \vec b$};
  \draw[attenzione] (0,0.22) -- (0.18,0.22) -- (0.18,0);
  % -c
  \draw[->, thick, attenzione!50, dashed] (0,0) -- (0,-1.1) node[below] {$\vec b \times \vec a = -\vec c$};
\end{tikzpicture}
